\section{Travaux réalisés : continuité du Projet de Fin d’Études}

La première phase des travaux a consisté en la conception et la validation d’un robot mobile autonome dans le cadre du Projet de Fin d’Études.

Ces travaux ont conduit à la publication de l’article :
\begin{quote}
“Design and Development of an Autonomous Mobile Robot for Unstructured Indoor Environments” \\
\cite{my_robot_paper}
\end{quote}

Ce papier présente :
\begin{itemize}
    \item la conception mécanique et électronique de la plateforme,
    \item l’intégration de capteurs (LiDAR, IMU, encodeurs),
    \item l’architecture logicielle sous ROS 2,
    \item la navigation autonome en environnement intérieur non structuré,
    \item la détection d’objets en temps réel.
\end{itemize}

Cette contribution constitue la plateforme expérimentale de référence pour la thèse.

À ce stade, les travaux portent principalement sur l’architecture robotique et l’autonomie navigationnelle. Aucun mécanisme d’apprentissage incrémental n’a encore été implémenté.
